\section*{Chapter 2: Background}
\setcounter{section}{2} \setcounter{subsection}{0}
\phantomsection \addcontentsline{toc}{section}{\textit{Chapter 2: Background}}

This chapter provides the technical context for the proposed framework. It reviews the principles of network slicing and the associated resource management models, discusses SLA-aware scheduling and two-tier control architectures, surveys the emerging body of work on LLM-based agents for network management, examines general-purpose and network-specific multi-agent coordination mechanisms, and concludes by identifying the research gaps that motivate the design presented in Chapter~3.

\subsection{Network Slicing Architecture and Resource Management Models}

Network slicing is one of the defining capabilities of 5G and beyond-5G networks. As surveyed by Afolabi \emph{et al.}~\cite{afolabi2018slicing}, slicing partitions a shared physical infrastructure into multiple end-to-end logical networks, each tuned to the requirements of a specific service class. The 3GPP service classification distinguishes three canonical slice types --- eMBB, URLLC and mMTC --- which differ substantially in their throughput, latency, reliability and connection-density targets. eMBB is designed for high-throughput applications such as video streaming and large file downloads, URLLC targets ultra-reliable and low-latency scenarios including industrial control and remote surgery, and mMTC supports massive concurrent device access typical of Internet-of-Things (IoT) sensor deployments. Because these slice types share a common pool of radio resources, resource allocation across slices must simultaneously respect hard capacity constraints and the heterogeneous performance objectives of each tenant.

Two broad families of methods have been applied to this problem. Rule-based and heuristic allocation schemes encode priorities and constraints explicitly and are attractive for their interpretability and bounded computational cost, but they generalise poorly to traffic patterns that deviate from those anticipated at design time. Learning-based approaches, and in particular DRL-based resource managers, have attracted considerable attention in recent years. Li \emph{et al.}~\cite{li2018drl} formulated bandwidth allocation across slices as a Markov Decision Process (MDP) and demonstrated that DRL agents can outperform hand-crafted heuristics under non-stationary demand. Sciancalepore \emph{et al.}~\cite{sciancalepore2019rlnsb} proposed the RL-NSB network slice broker, which uses reinforcement learning to admit and dimension slice requests so as to balance revenue and SLA compliance. Caballero \emph{et al.}~\cite{caballero2019games} analysed slicing under a game-theoretic lens, showing that slice customisation can be supported even when tenants behave strategically, and highlighting the inherently competitive nature of multi-tenant resource sharing. More recently, Nasser and Alani~\cite{nasser2025hybrid} explored hybrid deep-learning architectures for slicing, reinforcing the observation that learning-based methods can adapt to complex workload mixtures but require large amounts of training data and careful reward engineering.

A common limitation across these approaches is that the decision-making component is monolithic: a single policy is trained or hand-crafted to output allocations for all slices simultaneously. Neither the explicit articulation of the individual objectives of each slice nor a structured mechanism for reconciling conflicts between them is provided by the architecture itself. This limitation becomes particularly significant when slice demands change rapidly or when multiple SLA constraints must be balanced under tight resource budgets.

\subsection{SLA-Aware Scheduling and Two-Tier Control Mechanisms}

Slice resource management must operate simultaneously at two very different timescales. At the radio interface, scheduling decisions are taken on a per-Transmission Time Interval (TTI) basis --- typically every one or two milliseconds --- and must account for instantaneous channel quality, buffer occupancy and user fairness. At a longer timescale, inter-slice resource quotas, admission control decisions and policy updates are driven by aggregated performance indicators such as SLA satisfaction rates, tenant-level throughput and long-term utilisation.

This separation has led to the adoption of two-tier or hierarchical control architectures in which fast rule-based or model-based controllers handle per-TTI decisions, while slower optimisation or learning-based components adjust the parameters of the lower tier. Such architectures are compatible with the hierarchy already present in modern radio access networks, in which Medium Access Control (MAC)-layer schedulers operate at millisecond granularity while Radio Resource Management (RRM) functions revise resource partitions on a coarser timescale.

In the context of LLM-based decision making, the two-tier view is particularly relevant because the inference latency of contemporary LLMs is typically on the order of hundreds of milliseconds to several seconds per call, which renders it infeasible to place an LLM directly on the per-TTI control loop. Recent work on LLM-augmented O-RAN control~\cite{lotfi2025oranguide} and on dual-loop LLM-based management~\cite{qu2025dualloop} has therefore advocated the separation of a slow semantic reasoning loop from a fast execution loop, foreshadowing the two-tier design adopted in this project. The key engineering challenge lies in determining an appropriate decision frequency for the LLM layer, establishing a reliable fallback mechanism when inference latency exceeds acceptable thresholds, and ensuring that the lower tier can execute coherent allocation policies between successive upper-tier updates.

\subsection{LLM-Based Agent Frameworks in Network Management}

Between 2023 and 2025, research on LLMs for wireless networks developed along three broadly distinct technical routes, each making different trade-offs in how domain knowledge is injected into the model.

\subsubsection{Prompt Engineering and Agent Workflow Route}

The first route relies on prompt engineering and explicit agent workflows. WirelessAgent~\cite{tong2025wirelessagent} decomposes network management into a sequential workflow of intent understanding, slice allocation, bandwidth allocation and load balancing, where each stage is implemented as a LangGraph node that manages global state transitions. Its main contribution is demonstrating that a cognitive architecture consisting of perception, memory, planning and action, combined with tool calling, can approach rule-optimal performance in slicing scenarios without any task-specific training --- the reported bandwidth utilisation is only 4.3\% below the rule-optimal baseline under its considered experimental configuration, and 44.4\% higher than pure prompting methods. WirelessLLM~\cite{shao2024wirelessllm} proposes a complementary methodological roadmap: rapid domain knowledge injection through prompt engineering and Retrieval-Augmented Generation (RAG)~\cite{lewis2020rag} (knowledge alignment), fusion with specialised models (knowledge fusion), and model evolution through continual learning. Together, these works suggest that prompt-based approaches can be effective for network management when combined with appropriate architectural scaffolding, though they do not address the multi-agent coordination problem.

\subsubsection{Parameter Fine-Tuning Route}

The second route addresses domain adaptation through parameter fine-tuning. NetLLM~\cite{wu2024netllm} reformulates a selection of networking problems --- including viewport prediction, adaptive bitrate selection and cluster job scheduling --- as sequential decision problems and fine-tunes a Llama2-based backbone using low-rank adaptation (LoRA), reporting improvements over DRL baselines on all three tasks. While effective, this route requires substantial labelled data for each target task, a condition that is not easily satisfied in the network slicing domain where annotated operational datasets are scarce. Furthermore, fine-tuned models may lose their general reasoning capabilities, potentially limiting their applicability to novel or unforeseen network conditions.

\subsubsection{Cloud-Edge Collaboration Route}

The third route explores cloud-edge collaboration architectures. NetGPT~\cite{chen2024netgpt} proposes an architecture in which lightweight LLMs are deployed at the edge for latency-sensitive inference while larger models remain in the cloud for global optimisation, anticipating the deployment challenges posed by the size and inference cost of frontier models. Related work by Chergui \emph{et al.}~\cite{chergui2025collective} has investigated collective memory and knowledge sharing across LLM-based agents for cross-domain 6G management, exploring how distributed LLM agents can maintain coherent decision-making across heterogeneous network domains.

\subsubsection{Summary}

A comprehensive survey by Zhou \emph{et al.}~\cite{zhou2024llmtelecom} provides a broader overview of LLM applications across the telecommunications sector, documenting the rapid growth of this field and the diversity of approaches being explored. Across all three routes identified above, the decision-making agent is typically centralised: a single LLM or LLM-powered component is responsible for the entire resource management task. This observation underlies one of the research gaps identified later in this chapter.

\subsection{Multi-Agent Coordination and Negotiation Mechanisms}

\subsubsection{General-Purpose Multi-Agent LLM Frameworks}

In parallel with the network-specific work reviewed above, general-purpose multi-agent LLM systems have become an active area of research. MetaGPT~\cite{hong2023metagpt} encodes Standard Operating Procedures (SOPs) as prompt sequences and enforces the exchange of structured intermediate artefacts between agents, reducing the information loss that can arise when free-form natural language is the sole communication channel. AutoGen~\cite{wu2023autogen} provides a flexible conversational programming paradigm in which agents with different roles coordinate through multi-turn dialogues, and has been applied to tasks including code generation and data analysis. CAMEL~\cite{li2023camel} investigates the emergence of cooperative behaviour in role-playing agent societies, demonstrating that LLM agents can develop task-specific collaborative strategies through structured interaction. ChatDev~\cite{qian2023chatdev} applies multi-agent collaboration to software development by assigning dedicated roles to designers, coders and testers, showcasing the benefits of role specialisation in complex tasks. A survey by Wang \emph{et al.}~\cite{wang2023agentsurvey} provides a broader taxonomy of LLM-based autonomous agents and their cognitive components, identifying perception, planning, action and memory as the four foundational modules.

These frameworks have matured rapidly, but they share an important assumption: the underlying task is predominantly cooperative and admits a natural division of labour. Adapting them to network resource management raises three specific difficulties:

\begin{enumerate}
    \item Slice resource allocation is a \emph{competitive} constrained optimisation problem in which the total bandwidth is fixed and agents representing different slices must compete for a limited pool of resources.
    \item The decisions are subject to strong real-time constraints, which complicates the use of unbounded multi-turn dialogues.
    \item Decision outcomes are associated with clear quantitative evaluation criteria --- SLA satisfaction rates, throughput, latency --- which leave limited room for the qualitative judgement that dominates many general-purpose multi-agent benchmarks.
\end{enumerate}

\subsubsection{Multi-Agent Approaches in the Networking Domain}

Multi-agent attempts that target networking tasks directly remain limited. CommLLM~\cite{jiang2024commllm} proposes a three-component multi-agent architecture for 6G communications consisting of retrieval, cooperative planning and evaluation modules, but the reported validation is primarily focused on semantic communication and does not constitute a full slicing evaluation with quantitative resource allocation metrics. The dual-loop edge-terminal collaboration of Qu \emph{et al.}~\cite{qu2025dualloop} is architecturally informative in its treatment of an LLM-driven semantic loop coupled with a fast execution loop, yet coordination within that framework remains predominantly centralised, with limited inter-agent negotiation. ORAN-GUIDE~\cite{lotfi2025oranguide} adopts a division of labour in which an LLM provides semantic understanding and a reinforcement learning component handles online optimisation, offering useful insight into how LLM decision latency can be decoupled from network real-time requirements, but it does not implement multi-agent collaboration in the sense of multiple LLMs representing different network entities.

\subsubsection{Relevant Techniques}

Two techniques commonly adopted within these frameworks deserve specific mention. Chain-of-Thought (CoT) prompting~\cite{wei2022cot} encourages the model to reason step by step, and has been shown to improve performance on tasks involving arithmetic or logical decomposition; it is relevant for the structured resource allocation reasoning required of slice-manager agents. RAG~\cite{lewis2020rag} is frequently adopted as a lightweight mechanism for injecting domain knowledge without modifying model parameters, and can be used to provide agents with up-to-date network configuration data or SLA specifications.

\subsection{Identified Research Gaps}

Taken together, the work reviewed above suggests two complementary opportunities that this project seeks to address.

\textbf{Gap~1: Absence of an end-to-end multi-agent LLM framework for slice resource management.} Existing LLM-based network management frameworks either adopt a centralised single-agent design and do not expose explicit negotiation between slices~\cite{tong2025wirelessagent,shao2024wirelessllm}, or they propose multi-agent architectures whose validation is confined to communication tasks other than slice resource allocation~\cite{jiang2024commllm}. A complete framework dedicated to slice resource management that supports explicit multi-agent collaboration and is validated end-to-end in a reproducible slicing environment has not yet been presented in the accessible literature.

\textbf{Gap~2: Lack of structured negotiation protocols for competitive resource allocation.} Although general-purpose multi-agent LLM frameworks such as MetaGPT~\cite{hong2023metagpt} and ChatDev~\cite{qian2023chatdev} provide rich coordination mechanisms, these are designed with cooperative tasks in mind. Slice allocation is better modelled as a constrained multi-party optimisation problem in which the objectives of different slices compete under a fixed capacity budget~\cite{caballero2019games}. Structured negotiation and conflict-resolution protocols tailored to this competitive setting have received limited attention in the network management literature.

Table~\ref{table:comparison} summarises the positioning of this work relative to related frameworks across five key dimensions.

\begin{table}[h]
  \centering
  \begin{spacing}{1.35}
  \caption{Comparison of related frameworks and positioning of this research.}\vspace{0.7em}
  \label{table:comparison}
  \footnotesize
    \begin{tabular}{|p{2.6cm}|p{2.3cm}|p{2.3cm}|p{2.3cm}|p{3.5cm}|}
    \hline
    \textbf{Feature} & \textbf{WirelessAgent}~\cite{tong2025wirelessagent} & \textbf{CommLLM}~\cite{jiang2024commllm} & \textbf{Qu et al.}~\cite{qu2025dualloop} & \textbf{This Research} \\
    \hline
    Agent count & Single agent & Multi-agent (concept) & Dual-loop multi-agent & Multi-agent (3 slice + 1 coordinator) \\
    \hline
    Validation scenario & Network slicing & Semantic comm. & Edge-terminal collab. & Network slice resource allocation \\
    \hline
    Negotiation mechanism & None & Unspecified & Centralised coord. & Proposal--evaluation--negotiation protocol \\
    \hline
    Real-time handling & No explicit constraint & Not discussed & Inner-loop re-planning & Two-tier scheduling (LLM low-freq + rule high-freq) \\
    \hline
    Simulation validation & Custom env. & Concept level & Partial validation & Gymnasium standardised env. \\
    \hline
    \end{tabular}
  \end{spacing}
\end{table}

The design described in the following chapter responds to these two gaps by combining a two-tier scheduling architecture with a set of specialised slice-manager agents and a global coordinator that interact through an explicit proposal--evaluation--negotiation protocol.

%% ======== TEMPLATE GUIDANCE (retained per instructions) ========
% In this part of the report, you should give all the relevant background information about your project. Remember that your reader will not necessarily know the background technology you are using, so it is worthwhile to let them know.
%
% Also if your project is a research project, this is a good place to put down the related work or state of the art in the area – what the others have done? And why your research is novel?
%
% But don't make the background too long with every detail. It should be relevant to your project, with all the necessary information, written nicely and crisply.
